% =====================================================================
% Chương 6 — Lộ trình nhánh tiếp theo (Next Branch Roadmap)
% Nguồn thực thi: docs/roadmaps/next-branch.roadmap.json + .state.md
% =====================================================================

\chapter{Lộ trình nhánh tiếp theo}

Chương này trình bày lộ trình chi tiết cho nhánh phát triển tiếp theo của
dự án, xây dựng từ kết quả đối chiếu master plan hiện hành. Lộ trình được chia
thành \textbf{4 giai đoạn} ưu tiên theo đúng nguyên tắc ``đo được trước, sửa
lỗi ưu tiên cao, rồi mới thêm tính năng'':

\begin{enumerate}[label=\textbf{G\arabic*.}]
\item \textbf{G1 — Mở khả năng kiểm thử} (T1.1, T1.2, T1.3, T1.4): đưa các
thành phần phần mềm lên chạy/kiểm thử trên máy tính, không phụ thuộc board.
\item \textbf{G2 — Yêu cầu P0 còn lại} (T2.3, T3.4, T2.4): khôi phục cảnh báo
xe cắt ngang, thống nhất ngưỡng buzzer theo \code{FIRMWARE/SHARED}, nghiệm thu
link và dữ liệu cũ trên board.
\item \textbf{G3 — Tính năng đề xuất} (T3.1, T3.2, T3.3): biểu tượng vật thể,
sơ đồ xe tải EX8, mạch khuếch đại buzzer 5V.
\item \textbf{G4 — Vệ sinh \& đo lường} (T5.2, T5.3, T5.5--T5.9): nút
calibrate, comment sai, baseline, độ trễ, soak test.
\end{enumerate}

\section{Trạng thái kế thừa và nguyên tắc thiết kế}

Nhánh này kế thừa một số quyết định đã ghi nhận trong quá khứ, quyết định cả
cách phát triển:

\begin{center}
\small
\renewcommand{\arraystretch}{1.2}
\begin{longtable}{@{}p{6.6cm}p{9.8cm}@{}}
\caption{Trạng thái kế thừa ảnh hưởng lộ trình.}\\
\toprule
\textbf{Trạng thái} & \textbf{Ảnh hưởng tới nhánh này} \\
\midrule
\endhead
Mã nguồn đang ở \code{origin/main} (\code{6326682}), chưa có nhánh riêng
cho lộ trình. &
Tạo nhánh mới ngay từ commit này; mỗi giai đoạn một nhóm commit
conventional (R9). \\
Fix ESP-NOW (chủ động giữ kênh khi \code{WiFi.begin()} chạy song song,
tắt modem-sleep) \textbf{đã bị revert} theo yêu cầu. &
Không bóc tách sửa ``thử nghiệm''; đưa thành bước chính thức (G2) kèm
nghiệm thu board để tránh quay vòng revert. \\
\textbf{Buzzer:} \code{GPIO11} không thể xuất 5\,V và dòng đủ lớn —
cần transistor + rail 5\,V; \code{BUZZER\_PIN=11} (đã tránh GPIO47/48
xung đột PSRAM). &
G3.3 tài liệu mạch transistor, không cấp thẳng 5\,V từ GPIO. Cần người
dùng xác nhận buzzer \emph{active} hay \emph{passive}. \\
\textbf{CoreIoT:} sensor-node Inactive từ 14:59:36 (không publish MQTT). &
Xác nhận trước khi nghiệm thu G2: serial có \code{[NET] MQTT connected},
board đang flash env nào, telemetry \code{d1..d6} latest. \\
\textbf{Cảnh báo cắt ngang:} thuật toán cục bộ trong
\code{ui\_dashboard.c} chạy, nhưng nguồn ép từ rule-chain
(\code{crossing\_hazard}) là \emph{dead path} do node JS không xuất field. &
G2.5 gắn lại field \code{crossing\_hazard} đầu cuối, thống nhất macro
\code{CROSSING\_*} dùng chung trong \code{FIRMWARE/SHARED}. \\
Heuristic thực thi: \code{front\_close} \textless{} 150\,cm +
\code{side\_changing\_fast} (delta $\ge$ 40\,cm). &
Không đổi thông số cảnh báo, chỉ nối dead path và thêm test. \\
\bottomrule
\end{longtable}
\end{center}

\section{Toàn cảnh 15 bước và chuỗi phụ thuộc}

Chuỗi phụ thuộc cốt lõi: \textbf{G1 độc lập} → \textbf{G2 cần tìm board} →
\textbf{G3 tính năng cần simulator (G2.4)} → \textbf{G4 đo lường cần
record/replay (G2.3) và độ trễ (G4.14)}.

\begin{center}
\begin{tikzpicture}[
  node distance=5mm and 12mm,
  box/.style={draw, rounded corners=2pt, align=center,
              minimum height=7mm, font=\footnotesize, inner sep=1.5mm,
              fill=blue!5},
  pha/.style={font=\footnotesize\bfseries, draw=none, fill=none},
  arr/.style={-{Stealth[length=2.5mm]}, thick}
]
  % ---- Cột G1 ----
  \node[pha] (g1) {G1 — Kiểm thử};
  \node[box, below=1.5mm of g1] (s1) {1. T1.1 MQTT scenario};
  \node[box, below=2mm of s1] (s2) {2. T1.3a classify test};
  \node[box, below=2mm of s2] (s3) {3. T1.4 record/replay};
  \node[box, below=2mm of s3] (s4) {4. T1.2 LVGL SDL sim};
  % ---- Cột G2 ----
  \node[pha, right=18mm of g1] (g2) {G2 — P0 còn lại};
  \node[box, below=1.5mm of g2] (s5) {5. T2.3 crossing hazard};
  \node[box, below=2mm of s5] (s6) {6. T3.4 buzzer DANGER-only};
  \node[box, below=2mm of s6] (s7) {7. Nghiệm thu board};
  % ---- Cột G3 ----
  \node[pha, right=18mm of g2] (g3) {G3 — Tính năng};
  \node[box, below=1.5mm of g3] (s8) {8. T3.1 biểu tượng};
  \node[box, below=2mm of s8] (s9) {9. T3.2 sơ đồ EX8};
  \node[box, below=2mm of s9] (s10) {10. T3.3 buzzer 5V};
  % ---- Cột G4 ----
  \node[pha, right=18mm of g3] (g4) {G4 — Đo \& vệ sinh};
  \node[box, below=1.5mm of g4] (s11) {11. T5.2 calibrate};
  \node[box, below=2mm of s11] (s12) {12. T5.3 comment};
  \node[box, below=2mm of s12] (s13) {13. T5.5+5.6 baseline};
  \node[box, below=2mm of s13] (s14) {14. T5.7 latency};
  \node[box, below=2mm of s14] (s15) {15. T5.8+5.9 soak \& miss};

  \draw[arr] (s1.south) -- (s2.north);
  \draw[arr] (s2.south) -- (s3.north);
  \draw[arr] (s5.south) -- (s6.north);
  \draw[arr] (s6.south) -- (s7.north);
  \draw[arr] (s4.east) -- (s8.west);
  \draw[arr] (s7.east) -- (s11.west);
  \draw[arr] (s3.east) -- (s13.west);
  \draw[arr] (s13.south) -- (s14.north);
  \draw[arr] (s14.south) -- (s15.north);
\end{tikzpicture}
\captionof{figure}{Chuỗi phụ thuộc 15 bước của nhánh tiếp theo.}
\end{center}

\section{Bảng chi tiết từng bước}

Mỗi bước gồm nhóm file đích, định nghĩa hoàn thành (DoD) và lệnh xác minh
chạy được trên máy dev \textbf{không cần board} (trừ các bước ghi rõ).

\begin{center}
\footnotesize
\renewcommand{\arraystretch}{1.25}
\begin{longtable}{@{}p{3.4cm}p{4.4cm}p{8.6cm}@{}}
\caption{Khuôn mẫu DoD (R10) áp dụng mọi bước trong lộ trình.}\\
\toprule
\textbf{Thành phần} & \textbf{Nội dung} \\
\midrule
\endhead
\textbf{DoD phải đo được} &
Mỗi bước chỉ \textsc{Done} khi chạy được lệnh xác minh của chính bước đó
(\code{verification\_command}) và đầu ra như mô tả — \emph{``nhìn được''} là
không đủ. \\
\textbf{Kiểm tra lại hằng số} &
Sau mỗi bước sửa ngưỡng: chạy \code{check\_rulechain\_thresholds.py} và
\code{scan\_secrets.py} (R3/R11). \\
\textbf{Commit} &
Mỗi bước 1 nhóm commit conventional; cấm \code{git add -A}; chạy \code{/verify}
trước commit (R9). \\
\bottomrule
\end{longtable}
\end{center}

\subsection{G1 — Mở khả năng kiểm thử (bước 1--4)}

\begin{center}
\footnotesize
\renewcommand{\arraystretch}{1.3}
\begin{longtable}{@{}p{0.8cm}p{4.2cm}p{4.0cm}p{7.4cm}@{}}
\caption{G1 — các bước đưa phần mềm lên máy tính (không phụ thuộc board).}\\
\toprule
\textbf{\#} & \textbf{Nhiệm vụ} & \textbf{File đích} & \textbf{DoD (đo được)} \\
\midrule
\endhead
1 &
T1.1: \code{test\_mqtt\_coreiot.py} thêm \code{--scenario}
(\code{approach}/\code{crossing}/\code{slam}/\code{normal}) tự đặt chuỗi
khoảng cách theo thời gian, payload đúng schema V2
(\code{d1..d6} + \code{nearest\_cm} + \code{has\_nearest}). &
\code{tools/test\_mqtt\_coreiot.py}, \code{tools/guard/test\_guard.py} &
\code{python3 tools/test\_mqtt\_coreiot.py --scenario approach --dry-run
--distance 30} in 1 payload JSON đủ key; unit test guard pass. \\
2 &
T1.3: host test
\code{sensor\_model\_classify} tại ranh giới 30/100/0 (R3); thêm env
\code{native} cho waveshare. &
\code{firmware/waveshare-screen/test/test\_sensor\_model.cpp},
\code{platformio.ini} &
\code{pio test -e native} pass $\ge$ 6 case classify; không ảnh hưởng build
firmware thật. \\
3 &
T1.4: ghi \& phát lại telemetry JSONL — tái hiện cảnh không cần board. &
\code{tools/record\_telemetry.py}, \code{tools/replay\_telemetry.py} &
Ghi 10\,s ra \code{/tmp/tb.jsonl}, replay \code{--dry-run} không lỗi schema;
script có \code{--in/--out/--seconds/--topic}. \\
4 &
T1.2: simulator LVGL v9 + SDL2 render dashboard waveshare trên PC. &
\code{firmware/waveshare-screen/host\_sim/CMakeLists.txt},
\code{host\_sim/main.c}, \code{ui\_dashboard.h} &
\code{cmake -S ... -B /tmp/host\_sim \&\& cmake --build ...} thành công; chạy
được cửa sổ SDL (headless: xvfb) hoặc exit 0. \\
\bottomrule
\end{longtable}
\end{center}

\subsection{G2 — P0 còn lại (bước 5--7)}

\begin{center}
\footnotesize
\renewcommand{\arraystretch}{1.3}
\begin{longtable}{@{}p{0.8cm}p{4.2cm}p{4.0cm}p{7.4cm}@{}}
\caption{G2 — cảnh báo cắt ngang, thống nhất buzzer, nghiệm thu board.}\\
\toprule
\textbf{\#} & \textbf{Nhiệm vụ} & \textbf{File đích} & \textbf{DoD (đo được)} \\
\midrule
\endhead
5 &
T2.3: nối dead path \code{crossing\_hazard}. Node 3 rule-chain xuất thêm bool
JS (\code{FRONT $\le$ 150} \&\&\ \code{side delta $\ge$ 40}); firmware đọc tại
\code{main.c:84-87}; macro \code{CROSSING\_*} dùng chung qua
\code{firmware/shared}. &
\code{cloud/coreiot/rule\_chain/supersonic\_rule\_chain.json},
\code{firmware/waveshare-screen/components/ui\_dashboard/ui\_dashboard.c},
\code{firmware/waveshare-screen/src/main.c} &
\code{pio run -e yolo\_uno} OK; \code{grep crossing\_hazard .../rule\_chain/*.json}
hiện node JS; không còn phụ thuộc biến ép chết. \\
6 &
T3.4: thống nhất ngưỡng buzzer theo \code{FIRMWARE/SHARED} — bỏ
\code{BUZZER\_WARNING\_DISTANCE\_CM}/\code{BUZZER\_DANGER\_DISTANCE\_CM};
kêu khi \code{nearest\_cm $\le$ SENSOR\_DANGER\_CM} (30), im khi còn lại. &
\code{firmware/shared/thresholds.h},
\code{firmware/sensor-node/src/buzzer.cpp},
\code{firmware/sensor-node/test/test\_thresholds.cpp} &
2 env build OK; \code{pio test -e native} 10/10; grep 2 macro cũ = 0 (trừ
\code{.pio/}). \\
7 &
T2.4 + nghiệm thu: flash 2 board (sensor-node \code{yolo\_uno\_coreiot},
waveshare \code{yolo\_uno}), cùng AP. Xác nhận \code{LINKED} ổn định; slot
quá hạn về \code{--} cm; buzzer chỉ kêu $\le$ 30\,cm; crossing bật qua
rule-chain. &
\code{docs/PROGRESS.md}, \code{docs/logs/NEXT\_BRANCH\_ACCEPTANCE\_LOG.md} &
\emph{Flash-and-observe} ghi serial ngày-giờ; grep hàng DONE trong
\code{PROGRESS.md}. \textbf{Cần board + token thật.} \\
\bottomrule
\end{longtable}
\end{center}

\subsection{G3 — Tính năng đề xuất (bước 8--10)}

\begin{center}
\footnotesize
\renewcommand{\arraystretch}{1.3}
\begin{longtable}{@{}p{0.8cm}p{4.2cm}p{4.0cm}p{7.4cm}@{}}
\caption{G3 — các tính năng sau khi đã đo được (cần simulator bước 4).}\\
\toprule
\textbf{\#} & \textbf{Nhiệm vụ} & \textbf{File đích} & \textbf{DoD (đo được)} \\
\midrule
\endhead
8 &
T3.1: biểu tượng vật thể tại vị trí tương đối từng cảm biến trên canvas giữa
thay vì text. &
\code{firmware/waveshare-screen/components/ui\_dashboard/ui\_dashboard\_layout.c},
\code{ui\_dashboard.c} &
Simulator hiển thị icon di chuyển theo dữ liệu giả; firmware build OK; không
đổi log cảnh báo. \\
9 &
T3.2: sơ đồ xe tải \textbf{EX8} theo hồ sơ (hết ``hình gắn cứng''): vị trí 6
cảm biến + cung cảnh báo theo toạ độ profile. &
Như bước 8 + \code{firmware/shared/thresholds.h} &
Simulator hiển thị sơ đồ EX8; not chứa chuỗi \code{FRONT HOOD}/\code{REAR
TRUNK} trong mã nguồn (grep = 0 ngoài assets). \\
10 &
T3.3: mạch khuếch đại buzzer 5\,V qua transistor NPN/MOS —
\code{GPIO11} chỉ điều khiển (3.3\,V), buzzer ăn rail 5\,V; tài liệu sơ đồ
đấu dây. &
\code{firmware/shared/thresholds.h},
\code{docs/INSTALLATION\_SENSOR\_NODE.md},
\code{docs/HARDWARE\_INSTALLATION.md} &
2 env build OK; \code{INSTALLATION} có sơ đồ transistor (R 1\,k + pull-down
10\,k + rail 5\,V); bảng dây cập nhật. \emph{Cần trả lời: buzzer active/passive.} \\
\bottomrule
\end{longtable}
\end{center}

\subsection{G4 — Đo lường \& vệ sinh (bước 11--15)}

\begin{center}
\footnotesize
\renewcommand{\arraystretch}{1.3}
\begin{longtable}{@{}p{0.8cm}p{4.2cm}p{4.0cm}p{7.4cm}@{}}
\caption{G4 — các phép đo định lượng và vệ sinh cuối nhánh.}\\
\toprule
\textbf{\#} & \textbf{Nhiệm vụ} & \textbf{File đích} & \textbf{DoD (đo được)} \\
\midrule
\endhead
11 &
T5.2: nút \code{Calibrate} — hoặc gắn callback có ích (reset sensor về
nodata + log), hoặc bỏ nút. &
\code{ui\_dashboard\_layout.c}, \code{ui\_dashboard.c} &
Grep \code{calib\_btn} \code{add\_event\_cb} có (hoặc nút đã bỏ); build OK. \\
12 &
T5.3: đối chiếu comment ID cảm biến —
\code{sensor\_model.h} vs \code{espnow\_slot\_t}. &
\code{firmware/waveshare-screen/components/sensor\_model/include/sensor\_model.h},
\code{firmware/shared/espnow\_protocol.h} &
Grep 2 file mô tả cùng một thứ tự slot; label map không đổi thứ tự. \\
13 &
T5.5 + T5.6: đo baseline 6 khoảng cách (raw vs filter), tầm phản xạ ở 3 chiều
cao gá, búp thực tế vs \code{SENSOR\_BEAM\_FOV\_DEG} (75). &
\code{docs/PROGRESS.md}, \code{docs/TEST\_PROTOCOL.md} &
PROGRESS có bảng baseline $n=3$, ngày đo; kết quả cao/góc. \textbf{Cần board.} \\
14 &
T5.7: độ trễ đầu--cuối 2 nhánh (ESP-NOW vs MQTT) → \textbf{chốt T0.2} (nhánh
chính thức). &
\code{docs/PROGRESS.md} &
PROGRESS có bảng latency theo 2 nhánh + phương pháp đo timestamp (ESP-NOW
rssi\_frame, MQTT rule-chain). \textbf{Cần board.} \\
15 &
T5.8 + T5.9: soak 24/72\,h (heap, reset, link drop) và tỷ lệ báo nhầm/spot sót
trên tập replay đã gán nhãn (bước 3). &
\code{docs/PROGRESS.md} &
PROGRESS có mục soak (heap min, reset, uptime) + kết quả miss/false. \\
\bottomrule
\end{longtable}
\end{center}

\section{Rủi ro và điểm chặn chính}

\begin{itemize}
\item \textbf{Buzzer (bước 6, 10):} cần câu trả lời \emph{active/passive} và việc
lắp transistor trước khi tài liệu mạch 5\,V — nếu bỏ quên, chỉ thống nhất
ngưỡng phần mềm (bước 6), trễ mạch điện (bước 10).
\item \textbf{Bước 7 (nghiệm thu):} phụ thuộc board vật lý + token thật từ
\code{config/keys.json} (R1). Vấn đề ESP-NOW đã revert cần nghiệm thu trên hiện
trường, không quay vòng sửa-đoán.
\item \textbf{Bước 4 (simulator):} rủi ro phiên bản LVGL/SDL trong
\code{managed\_components}; không gộp simulator vào CMake firmware thật, giữ
thành project riêng \code{host\_sim/}.
\item \textbf{Bước 13--15 (đo):} tốn board + không gian; có thể hoãn 13--15
sang nhánh sau nếu tài nguyên không cho phép — luôn giữ G1 (1--4) làm tiền đề
đo mới mà không cần board.
\item \textbf{Quy tắc tập trung:} tối thiểu thay đổi ngoài phạm vi — mỗi bước
chỉ sửa đúng file đích đã khai; tránh ``sửa lan'' gây revert (bài học từ fix
ESP-NOW).
\end{itemize}